\section{Conclusion}
\label{sec:conclusion}

This paper presents QoS-FiLM, a graph neural network architecture that uses per-flow QoS intent as a conditioning signal to address online routing under heterogeneous traffic. QoS-FiLM injects flow intent into every message-passing layer of the GNN through feature-wise linear modulation, so that the same physical topology produces differentiated node embeddings for flows with different QoS requirements, resolving the ``same graph, same embedding'' problem of representational homogeneity. Experiments show that QoS-FiLM outperforms classical baselines in cumulative utility, bandwidth satisfaction ratio, and end-to-end latency across multiple backbone topologies. Ablation studies confirm that per-layer FiLM conditioning yields a performance gain of $+0.219$ over the flow-unaware variant, that the choice of graph attention operator strongly modulates this gain, and that FiLM and the attention operator interact synergistically rather than additively. Interpretability analysis verifies that the FiLM encoder is highly sensitive to QoS intent at the representation level while the flow-unaware baseline cannot perceive intent changes. 
Future work will (i) extend the approach to multi-domain topologies, and adaptive candidate-path generation, and (ii) investigate meta-learning strategies for rapid adaptation to unseen network topologies.
